\honest{An Intuitive Explanation of Bayes's Theorem}{Eliezer Yudkowsky}{2003}

A note at the top says that I now consider this explanation obsolete. \nb{It is still the version served under this title.}

Your friends are excited about Bayes's Theorem, and you look it up and find a single equation. ``What is the so-called Bayesian Revolution now sweeping through the sciences''? ``What is the light that they have seen?'' And then: ``Soon you will know. Soon you will be one of us.''

Experts say that Bayesian reasoning is hard for humans to grasp: people do not use it by instinct, learn it with difficulty and forget it after tutoring. ``Or so they claim.'' I promise ``an excruciatingly gentle introduction,'' using natural frequencies and pictures, to show what the numbers mean and why the rules must be what they are.

The problem is this. Of women screened at forty, 1\% have breast cancer. Of those, 80\% get a positive mammogram, and 9.6\% of healthy women get one too. A woman tests positive. What is the chance she has cancer? ``Only around 15\% of doctors'' get it right, and I cite three studies. \nb{None of the three reports that figure for doctors on this problem; the study whose numbers come closest tested students in Salzburg.} Is the figure an urban legend? No: ``It's a surprising result which is easy to replicate, so it's been extensively replicated.'' I name no study beyond the three. Doctors do better when the same problem is stated as counts, 10 in 1,000 and so on, and best, 46\% correct, when it is stated as 100 of 10,000 women, 80 of the 100, and 950 of the 9,900.

Then the method that works. Imagine 10,000 women. A hundred have cancer, and 80 of them test positive. Of the 9,900 without cancer, 950 test positive. So 80 of the 1,030 women who test positive have cancer: 7.8\%, roughly one in thirteen. \nb{It is Gigerenzer and Hoffrage's method of counting people instead of multiplying probabilities, and it is the best thing in the essay.} After the test the women fall into four groups: 80 with cancer and a positive result, 20 with cancer and a negative one, 950 healthy and positive, 8,950 healthy and negative.

Most doctors answer between 70\% and 80\%, which is ``wildly incorrect.'' They replace the 1\% with the 80\%, ignoring both how rare the cancer is and how often healthy women test positive. The answer needs all three numbers. The first, 1\%, is the prior probability. The two rates of positive results, for sick and for healthy women, are the conditional probabilities. The answer, 7.8\%, is the posterior probability.

Where only one woman in a million has cancer, a test that catches 8 cancers in 10 and gives false positives 1 time in 10 still yields a hundred thousand false alarms per real case. The result raises the odds from 1 in a million only to 1 in 100,000. A result does not replace what you knew; it slides the probability up or down. The chance that a woman who tests positive has cancer and the chance that a woman with cancer tests positive are ``as unlike as apples and cheese.''

A second example. A barrel holds plastic eggs; 40\% contain pearls. Of the pearl eggs, 30\% are painted blue. Of the empty ones, 10\% are. What share of blue eggs hold pearls? Of all eggs, 12\% are blue with pearls and 6\% blue and empty, so 12 of 18 blue eggs, about 67\%, hold pearls. With pearls in one egg in a thousand, blue would raise the chance only from 0.1\% to 0.3\%; with pearls in 999 of 1,000, from 99.9\% to 99.966\%.

People new to this answer 30\%, or 20\% by subtracting the false positive rate from the true one. I tell of second-graders who answer ``how old is the bus driver?'' by adding up the passengers: they know which procedure is being prompted but have not connected it to reality. I give no source for that story.

Gigerenzer and Hoffrage found that stating the problem in frequencies works slightly better than stating it in probabilities, and that counts work best of all. \nb{Their study found no gain for the middle step.} The best format, natural frequencies, builds the prior into the conditional numbers: of 100 eggs, 40 hold pearls and 12 of those are blue, while 60 are empty and 6 of those are blue. It is what you would see cracking eggs yourself. Even so, only about half of people then reason correctly, not enough for real doctors and patients.

Where do priors come from? ``Never ask that question.'' Scientists' priors are set ``by annual vote of the AAAS,'' and everyone else downloads theirs ``from Kazaa.'' Then I give a serious answer. Priors are true or false like any answer, and can be checked against reality, for instance by counting how many women in a sample have cancer. Here three studies would supply the three numbers. \nb{That works for the rate of a disease. It does not work for the prior probability of a scientific theory, which has no population to count.}

In the formal part, the joint probability $P(A,B)$ equals $P(B,A)$, but $P(A|B)$ is not $P(B|A)$, and neither is $P(A,B)$. I count degrees of freedom. $P(\mathrm{cancer})$ and $P(\neg\mathrm{cancer})$ have one between them, since they sum to 1. The two positive rates, for sick and for healthy women, have two, since either can be anything. $P(\mathrm{positive},\mathrm{cancer}) = P(\mathrm{positive}|\mathrm{cancer}) \times P(\mathrm{cancer})$ ties three quantities together. Tests are calibrated this way: dividing, say, 6,816 positives by 8,520 women with cancer gives 80\%. Dividing the other way gives nonsense like 125\%, a common slip. The four groups have three degrees of freedom, since their fractions sum to 1, so any three independent numbers, such as the prior and the two conditional probabilities, fix the whole problem.

The ratio of the true positive rate to the false positive rate is the likelihood ratio, which says how far a positive result slides the probability. It does not say what a negative result means, or how often the test helps. A test with 80\% hits and 9.6\% false positives has the same ratio as one with 8\% and 0.96\%, but the first is better in every way.

Tests can be combined if they are independent. Add an invented second test, with 90\% hits and 5\% false positives, and a woman positive on both has about a 60\% chance of cancer. Odds make the bookkeeping easy: a 1\% prior is odds of 1 to 99, and three independent tests with likelihood ratios 25:3, 18:1 and 7:2 multiply to 3,150 to 594, or 84\%. The order of the tests does not matter: ``The proof is left as an exercise for the reader.''

E. T. Jaynes suggested measuring evidence in decibels, ten times the base-10 logarithm, so that evidence adds: the prior is about $-20$ decibels, the tests add 9, 13 and 5, and the result, 7 decibels, is odds of about 5 to 1, around 83\%.

A gizmo problem, blocked hoses and sparks, shows the arithmetic in general: multiply the prior by the chance of the evidence if the hypothesis is true, and divide by that plus the same product for the alternative. In general form, $P(A|X) = P(X|A)P(A) / [P(X|A)P(A) + P(X|\neg A)P(\neg A)]$, for a phenomenon A and evidence X. Late in the essay I say that if Bayes's Theorem now seems obvious, ``this introduction has entirely succeeded in its purpose.''

Bayes's Theorem says what counts as evidence and how much. The Bayesian method, I say, sets the limit on what a piece of evidence can yield, as thermodynamics limits the work from a temperature difference, and in cognitive science ``Bayesian reasoner'' is the precise word for a rational mind.

People neglect prior frequencies. They also attend to how likely the evidence is if A is true, and not to how likely it is if A is false. Rain nearly guarantees wet grass, but wet grass does not prove rain, since sprinklers and dew wet it too. If grass were never wet without rain, wet grass would prove rain even if rain wet it only half the time. ``Evidence is always the result of the differential between the two conditional probabilities.'' And: ``Strong evidence is not the product of a very high probability that A leads to X, but the product of a very low probability that not-A could have led to X.''

As for the question I opened with, the Bayesian revolution is fueled by ``more and more cognitive scientists'' and by ``scientists in every field,'' none of whom I name. Science itself is a special case of Bayes's Theorem: whether an experiment confirms a theory depends also on whether other explanations predict the same result. Karl Popper held that theories can be ``definitely falsified.'' \nb{Popper also wrote that in practice no conclusive disproof of a theory can ever be produced.} On the Bayesian account, a theory can make the evidence as likely as it likes, but it cannot control how likely rival theories make it, so confirmation has a ceiling. Newton's theory of gravity was toppled by ``the hidden gotcha,'' evidence that a rival theory predicted and Newton's did not. \nb{That evidence had been in plain view for decades: the anomaly in Mercury's orbit was known from 1859. Newton's theory fell only in 1915, when Einstein's theory explained it.}

Evidence a theory says is nearly impossible counts hugely against it: a result with probability 0.0001\% under the theory and 1\% otherwise is a ratio of 1 to 10,000, $-40$ decibels. Falsifiability follows from conservation of probability: if X would confirm a theory, not-X must count against it. But Popper was wrong that there is no confirmation at all. Falsification is only much stronger, and it is still probabilistic.

Then I go back to the notation and explain it again. A better notation exists, ``which it is now too late to adopt.'' In $P(A|X)$, A is what you want to know and X the evidence. $P(Q,P)$ is the share of all things that have both properties, say 641 of 89,031 women, 0.72\%. By contrast, $P(Q|P)$ is the share among those that have P, 641 of the 7,915 who tested positive. Conditioning on P shrinks your world to the things with P. The right side of the theorem follows from the left in a few steps: $P(A|X) = P(X,A)/P(X)$, and $P(X)$ is $P(X,A) + P(X,\neg A)$.

Finally: ``Rational inference on the left end, physical causality on the right end; an equation with mind on one side and reality on the other.'' \nb{The theorem holds whatever causes what. Two sentences earlier Yudkowsky had said only that causal relations ``generally'' run this way.} ``You are now an initiate of the Bayesian Conspiracy.''
